\section*{Reading rules and status legend}

\paragraph{Objects.} A polygon is a cyclic orbit of labelled tuples of points
of $\RR^2$ (Definition~\ref{def:polygon}); every construction is made on a
representative with labels $1,\ldots,n$ and is shift-equivariant, and a root
or base vertex, when one is needed, is part of the construction's data.
``Generic'' has one meaning (Definition~\ref{def:generic}). Each wall type is a named germ
(Definition~\ref{def:walls}); no law is asserted at a configuration that is not
a named germ.

\ifprintstatus
\paragraph{Status tags.} Every definition-free statement ends with one tag.
A tag records the provenance and the audit state of the printed proof; it is
a claim of this document, and the audit ledger of Phase~2.5 is the status.
\begin{itemize}
\item \textsc{proved}: a complete proof is printed here and either carried
unchanged from frame SM1 --- such items were carried from SM1 unrefereed, read
by both benches on SM12, repaired in rounds~12--13 and re-read on SM13 and SM14;
all eight carry the refereed tag from frame SM15 (CUR-PROC-4) --- or printed in round~1 from a
frozen campaign document and refereed on frame SM2 by the auditor bench~A on
both axes, fidelity to the source and the argument (ruling R-25-20); in the
second case the tag reads \textsc{proved (refereed: bench A, stamp;
transcribed from locator)} and keeps the locator; an item cleared again
after a directed touch lists both stamps, and Lemma~\ref{lem:carriers}
carries the tag on its clauses (i)--(iii) with clause (iv) named as a new
argument. The tag remains a claim of
this frame; the audit ledger is the status.
\item \textsc{transcribed (unrefereed)}: a complete proof is printed here
(round~1 of Phase~2.5, or so retagged in round~2), written from the frozen
campaign document named after the colon (CV = Phase~1.999 output, frame CV,
files \texttt{d1}--\texttt{d10}; RC = Phase~1.98 release candidate RC\_v4; RA =
the R3 attachment frame; appendix\_v7 = the main text's appendix file; SM1 =
frame SM1 of this phase, for a proof whose SM1 bytes were paraphrased in
round~1) in the present conventions;
neither its fidelity to that source nor the argument itself has been refereed
yet. A held transcription whose statement was repaired in round~3, round~4,
round~5, round~6, round~8 or round~9 keeps this tag with the repair named after the locator; one whose proof
text was changed in round~3, round~4, round~5, round~6, round~8 or round~9 is tagged \textsc{new} with the locator
kept. A cleared item touched by a directed repair reverts to this tag with
its clearance history recorded, until re-refereed. The transcription table is Section~\ref{sec:transcriptions}. No statement of this frame carries this tag.
\item \textsc{open}: a proof obligation not discharged anywhere. No statement
of this frame carries this tag; the open obligations are listed in
Section~\ref{sec:open}.
\end{itemize}
\fi

{\raggedright
\paragraph{Frozen sources.} The transcription sources are CV,
\path{phase199/manuscript/frames/CV/}; RC,
\path{phase198/manuscript/RC_v4/}; and RA,
\path{phase1999/evidence/codex/w8_ra_frame_v1/RA/}.
The appendix source is the curator-posted read-only copy of appendix v7.
Locators name labels in those documents. CV uses the same one-based tail
labels as this document. RC uses zero-based head labels and
$[u,v]_*=-\det(u,v)$; its physical endpoint map is stated in
Section~\ref{sec:conventions}. No starred bracket is used here.\par}
